% Part: second-order-logic
% Chapter: syntax-and-semantics
% Section: semantic-notions

\documentclass[../../../include/open-logic-section]{subfiles}

\begin{document}

\olfileid{sol}{syn}{sem}
\olsection{અર્થવિચારી ખ્યાલો}

\begin{explain}
\emph{પ્રામાણ્ય}, \emph{ફલિતતા} અને \emph{સંતોષ્યતા}ના કેન્દ્રીય
તાર્કિક ખ્યાલો દ્વિતીય-ક્રમ તર્કશાસ્ત્રમાં પ્રથમ-ક્રમ તર્કશાસ્ત્રની
જેમ જ વ્યાખ્યાયિત થાય છે; માત્ર તેમનો પાયાનો સંતોષ-સંબંધ હવે
દ્વિતીય-ક્રમ !!{formula}s માટેનો છે. અલબત્ત, દ્વિતીય-ક્રમ
!!{sentence} એવું !!a{formula} છે જેમાં વિધેયધર્મ-ચલો અને
વિધેય-ચલો સહિત બધા ચલો બદ્ધ હોય છે.
\end{explain}

\begin{defn}[પ્રામાણ્ય]
વાક્ય~$!A$ \emph{પ્રમાણભૂત}, $\Entails !A$, ત્યારે અને માત્ર ત્યારે જ
જ્યારે દરેક !!{structure}~$\Struct M$ માટે $\Sat{M}{!A}$.
\end{defn}

\begin{defn}[ફલિતતા]
વાક્યોનો ગણ~$\Gamma$, વાક્ય~$!A$ને \emph{ફલિત કરે} છે,
$\Gamma \Entails !A$, ત્યારે અને માત્ર ત્યારે જ જ્યારે
$\Sat{M}{\Gamma}$ હોય એવી દરેક !!{structure}~$\Struct M$ માટે
$\Sat{M}{!A}$.
\end{defn}

\begin{defn}[સંતોષ્યતા]
વાક્યોનો ગણ~$\Gamma$ \emph{સંતોષ્ય} છે જો કોઈ !!{structure}~$\Struct M$
માટે $\Sat{M}{\Gamma}$. જો~$\Gamma$ સંતોષ્ય ન હોય, તો તેને
\emph{અસંતોષ્ય} કહે છે.
\end{defn}

\end{document}
